%%%%%%%%%%%%%%%%%%%%%%%%%%%%%%%%%%%%%%%%%%%%%%%%%%%%%%%%%%%%%%%%%%%%%%%%
%                                                                      %
%     File: Thesis_Abstract.tex                                        %
%     Tex Master: Thesis.tex                                           %
%                                                                      %
%     Author: Andre C. Marta                                           %
%     Last modified :  4 Mar 2024                                      %
%                                                                      %
%%%%%%%%%%%%%%%%%%%%%%%%%%%%%%%%%%%%%%%%%%%%%%%%%%%%%%%%%%%%%%%%%%%%%%%%

\section*{Abstract}

% Add entry in the table of contents as section
\addcontentsline{toc}{section}{Abstract}

The use of prosthetic devices in animals has increased significantly in recent years, driven by advances in veterinary medicine and biomedical engineering. However, the evaluation of prosthesis acceptance and effectiveness remains largely based on qualitative observation, relying on subjective assessment by veterinarians, which limits consistency and reproducibility.

This work proposes the development of an instrumented prosthesis to monitor canine locomotion through the acquisition of relevant biomechanical data. The system integrates inertial and force sensors to measure parameters such as acceleration, orientation, and load patterns during gait, using a low-power embedded system with local data storage.

Since sensing is limited to the prosthetic limb, adaptation cannot be assessed using traditional inter-limb symmetry metrics. Instead, proxy metrics based on temporal consistency, signal variability, and force distribution patterns are proposed to infer the level of adaptation.

Based on these data, machine learning models will be explored to automatically classify the adaptation state of the prosthesis, framing the problem as a supervised classification task. Model selection will be addressed in a later stage, according to the characteristics of the acquired data.

This approach aims to support the transition from qualitative to quantitative and automated evaluation methods, improving the adaptation and optimization of prosthetic devices. Furthermore, the proposed methodology may have broader implications for human prosthetics, contributing to more reliable, data-driven evaluation techniques.

\vfill

\textbf{\Large Keywords:} canine prosthetics, gait analysis, embedded systems, sensors machine, inertial  learning



